% !TeX root = ../../main.tex
\subsection{\bilingual{Coupling of \texorpdfstring{$N$}{N}-BRW and Killed BRW}{\texorpdfstring{$N$}{N}-分支随机游走与经灭杀的分支随机游走的耦合}}
\begin{lemma}
    \label{lemma:4.8}
    \bilingual{
        Let $\mathcal V^{\mathrm k}$ be the branching random walk $\mathcal V$ killed according to a given rule, such that the generation size process $\#\mathcal Z^{\mathrm k}$ is adapted to the natural filtration generated by $\mathcal V$. Assume that at time $0$, the population of $\mathcal V^{\mathrm k}$ is included in (or equal to) that of $\mathcal V^N$. Then there exists a coupling of $\mathcal V^N$ and $\mathcal V^{\mathrm k}$ such that
    }{
        设 $\mathcal V^{\mathrm k}$ 是按给定灭杀规则经灭杀的分支随机游走,~ 使得其粒子总数过程 $\#\mathcal Z^{\mathrm k}$ 适应于 $\mathcal V$ 的自然流.~ 假设在时刻 $0$,~ $\mathcal V^{\mathrm k}$ 的粒子分布包含于 (或等于) $\mathcal V^N$ 的粒子分布.~ 则存在 $\mathcal V^N$ 和 $\mathcal V^{\mathrm k}$ 的一个耦合使得
    }
    {
        \setlength{\belowdisplayskip}{-10pt}
        \begin{equation*}
            \mathcal M^N_k
            \le
            \mathcal M^{\mathrm k}_k,
            \quad
            k\le\tau:=\inf\left\{
                k\ge 0:
                \#\mathcal Z^{\mathrm k}_k
                >
                N
            \right\}.
        \end{equation*}
    }
\end{lemma}

\begin{proof}
    \quad

    \bilingual{
        For two finite counting measures on $\mathbb R$, i.e.,
    }{
        对于 $\mathbb R$ 上的两个有限计数测度,~ 即
    }
    $
        \mu=\sum_{i=1}^{a(\mu)}\delta_{x_i}$,
        % \quad
        $\nu=\sum_{i=1}^{a(\nu)}\delta_{y_i}$,
    % \end{equation*}
    \bilingual{
        where $x_1\le x_2\le\cdots\le x_{a(\mu)}$ and $y_1\le y_2\le\cdots\le y_{a(\nu)}$. It suffices to construct a coupling of $\mathcal V^N$ and $\mathcal V^{\mathrm k}$ such that
    }{
        其中 $x_1\le x_2\le\cdots\le x_{a(\mu)}$ 且 $y_1\le y_2\le\cdots\le y_{a(\nu)}$.~ 只需构建 $\mathcal V^N$ 和 $\mathcal V^{\mathrm k}$ 的一个耦合使得
    }
    \begin{equation}
        \mathcal V^{\mathrm k}_k
        \succcurlyeq
        \mathcal V^N_k,\quad
        k\le\tau.
        \label{eq:4.12}
    \end{equation}
    \bilingual{
        We proceed by induction on $k$. By assumption, \eqref{eq:4.12} holds for $k=0$. Assume that the coupling \eqref{eq:4.12} has been built up to some time $k<\tau$. We may write $\mathcal V^{\mathrm k}_k=\sum_{i=1}^{a_k}\delta_{x_i}$ and $\mathcal V^N_k=\sum_{i=1}^{b_k}\delta_{y_i}$ such that $a_k\le b_k$ and $x_i\ge y_i$ for all $i\in[\![1,a_k]\!]$. If $a_k=0$, $\tau=\infty$, and \eqref{eq:4.12} holds for all times after $k$ by definition. Consider the non-degenerate case where $a_k\ge 1$. We construct the coupling at time $k+1$ as follows.
    }{
        对 $k$ 归纳证明.~ 根据假设,~ \eqref{eq:4.12} 在 $k=0$ 时成立.~ 假设耦合 \eqref{eq:4.12} 已经构建到某个时刻 $k<\tau$.~ 此时可以写 $\mathcal V^{\mathrm k}_k=\sum_{i=1}^{a_k}\delta_{x_i}$ 和 $\mathcal V^N_k=\sum_{i=1}^{b_k}\delta_{y_i}$ 使得 $a_k\le b_k$ 且对于所有 $i\in[\![1,a_k]\!]$,~ 都有 $x_i\ge y_i$.~ 如果 $a_k=0$,~ 则 $\tau=\infty$,~ 根据定义 \eqref{eq:4.12} 在 $k$ 之后的所有时刻都成立.~ 考虑非退化情况 $a_k\ge 1$.~ 我们如下构建时刻 $k+1$ 的耦合.
    }
    \begin{itemize}
        \item
        \bilingual{
            For $1\le i\le a_k$, each $x_i$ particle of $\mathcal V^{\mathrm k}_k$ produces a collection of offspring $x_i+\Xi_{i,k}$ at time $k+1$ in $\mathcal V$, where $(\Xi_{i,k})_{i\ge 1,k\ge 0}$ are i.i.d.~copies of $\Xi$. Those offspring that are non-killed according to the given rule (if the set is not empty) constitute $\mathcal V^{\mathrm k}_{k+1}$. Note that if the killing cannot be performed here due to the use of future information in the killing rule, and the number of offspring is not well defined, then the proof cannot proceed, which is exactly the situation excluded by this lemma.
        }{
            对 $1\le i\le a_k$,~ $\mathcal V^{\mathrm k}_k$ 的每个粒子 $x_i$ 在时刻 $k+1$ 产生一组后代 $x_i+\Xi_{i,k}$,~ 其中 $(\Xi_{i,k})_{i\ge 1,k\ge 0}$ 是 $\Xi$ 的独立同分布的拷贝. 根据给定的灭杀规则,~ 那些未被灭杀的后代 (如果该集合非空) 构成 $\mathcal V^{\mathrm k}_{k+1}$. ~注意,~ 如果这里由于灭杀规则用到了未来的信息而无法执行灭杀,~ 从而无法确定余下粒子的数量,~ 这里的证明就无法继续下去,~ 这也就是本引理所排除的情形.
        }
        \item
        \bilingual{
            For $\mathcal V^N$, each $y_i$, $i\in[\![1,b_k]\!]$, produces a collection of offspring $y_i+\Xi_{i,k}$ at time $k+1$, and we select the $N$ leftmost particles to form $\mathcal V^N_{k+1}$. In particular, $\#\mathcal V^N_{k+1}=N$.
        }{
            对 $\mathcal V^N$,~ 每个 $y_i$,~ $i\in[\![1,b_k]\!]$,~ 在时刻 $k+1$ 产生一组后代 $y_i+\Xi_{i,k}$,~ 选取最左侧的 $N$ 个粒子组成 $\mathcal V^N_{k+1}$.~ 特别地,~ $\#\mathcal V^N_{k+1}=N$.
        }
        \item
        \bilingual{
            If the population of $\mathcal V^{\mathrm k}_{k+1}$ does not exceed $N$, then necessarily $\mathcal V^{\mathrm k}_{k+1}\succcurlyeq\mathcal V^N_{k+1}$, and \eqref{eq:4.12} holds at time $k+1$. Otherwise, we have $\#\mathcal V^{\mathrm k}_{k+1}>N$, which implies that $\tau=k+1$ and \eqref{eq:4.12} holds as well.
        }{
            若 $\mathcal V^{\mathrm k}_{k+1}$ 的粒子总数不超过 $N$,~ 则必然 $\mathcal V^{\mathrm k}_{k+1}\succcurlyeq\mathcal V^N_{k+1}$,~ 于是 \eqref{eq:4.12} 在时刻 $k+1$ 成立.~ 否则,~ 有 $\#\mathcal V^{\mathrm k}_{k+1}>N$,~ 此时 $\tau=k+1$ 且 \eqref{eq:4.12} 同样成立.
        }
        \qedhere
    \end{itemize}
\end{proof}

\begin{figure}[H]
    \centering
    \begin{tikzpicture}[
        xscale=3.3,
        yscale=1.1,
        >=Stealth,
        % 样式定义
        node-both/.style={circle, fill=blue!70!black, draw=green!60!black, line width=1.2pt, inner sep=1.2pt},
        node-N/.style={circle, fill=blue!70!black, inner sep=1.8pt},
        node-k/.style={circle, fill=green!60!black, inner sep=1.8pt},
        node-N-dead/.style={circle, draw=blue!70!black, fill=white, inner sep=1.4pt, cross out, draw=blue!80!black, line width=1.2pt},
        node-k-dead/.style={circle, draw=green!60!black, fill=white, inner sep=1.4pt, cross out, draw=green!60!black, line width=1.2pt},
        node-dead/.style={circle, draw=black!40, fill=white, inner sep=1.4pt},
        % 连线样式
        edge-N/.style={thick, blue!70!black},
        edge-k/.style={thick, green!60!black},
        edge-N-dashed/.style={dashed, blue!70!black},
        edge-k-dashed/.style={dashed, green!60!black},
        edge-both/.style={
            thick,
            blue!70!black, % 底色：蓝色
            postaction={
                draw,
                green!60!black,
                % on 2pt 代表绿色段长度，off 2pt 代表空隙（露出底下的蓝色）
                % 若要更频繁，将 2pt 改为 1pt
                dash pattern=on 5pt off 5pt
            }
        },
        edge-dashed/.style={dashed, black!30},
        % 灭杀边界
        kill-line/.style={red!80!black, dotted, ultra thick}
    ]

        % --- 坐标轴与背景 ---
        % \draw[->, gray!50, thick] (-0.2, -4.5) -- (-0.2, 2.5);
        \foreach \x in {0,1,2,3} \draw[-, black!20, dashed] (\x, -4.5) -- (\x, 2.5);
        \node[anchor=south] at (0, 2.7) {$k=0$};
        \node[anchor=south] at (1, 2.7) {$k=1$};
        \node[anchor=south] at (2, 2.7) {$k=2$};
        \node[anchor=south] at (3, 2.7) {$\tau$};

        % 灭杀屏障 (Vk 专用)
        % \draw[kill-line] (0, 0.2) -- (1, 0.2);

        % ================= n=0 =================
        \node[node-both] (n0) at (0, 0) {};
        \node[left] at (n0) {{\color{black}$\Xi_{1,0}$}};

        % ================= n=1 (N=2) =================
        % 产生4个后代。VN选前2个，Vk选1,3名（第2名过界被杀）
        \node[node-both] (n1-1) at (1, -1.2) {};
        \node[below] at (n1-1) {{\color{black}$\Xi_{1,1}$}};
        \node[node-N]    (n1-2) at (1, -0.2) {}; % VN保留，但Vk在此处被灭杀（过界或规则）
        \node[below] at (n1-2) {{\color{blue!70!black}$\Xi_{2,1}$}};
        \node[node-k]    (n1-3) at (1, 0.4)  {}; % VN淘汰（排名第3），但Vk保留
        \node[above] at (n1-3) {{\color{green!60!black}$\Xi_{2,1}$}};
        \node[node-N-dead] (n1-4) at (1, 1.5)  {};

        \draw[edge-both] (n0) -- (n1-1);
        \draw[edge-N] (n0) -- (n1-2); \draw[edge-N-dashed] (n0) -- (n1-2);
        \draw[edge-N-dashed] (n0) -- (n1-3); \draw[edge-k] (n0) -- (n1-3);
        \draw[edge-N-dashed] (n0) -- (n1-4);

        % ================= n=2 (N=2) =================
        % 耦合发生：VN从 y1, y2 繁殖；Vk从 x1, x3 繁殖

        % 1. VN 的前沿从 n1-2 爆发（产生极低值）
        \node[node-N] (n2-N-min) at (2, -3.0) {};
        \draw[edge-N] (n1-2) -- (n2-N-min);
        \draw[edge-N] (n1-2) -- (2, -2.2);
        \node[node-N] at (2, -2.2) {};
        \draw[edge-N-dashed] (n1-2) -- (2, -0.4);
        \node[node-N-dead] at (2, -0.4) {};

        % 2. Vk 的前沿从 n1-1 延续
        \node[node-k] (n2-k-min) at (2, -0.6) {};
        \node[below] at (n2-k-min) {{\color{green!60!black}$\Xi_{1,2}$}};
        \node[node-N-dead] at (2, -0.6) {};
        \draw[edge-both] (n1-1) -- (n2-k-min);

        % 3. 展现绿色来自被蓝色去掉的祖先
        \node[node-k] (n2-k-from-dead) at (2, 0.2) {};
        \node[above] at (n2-k-from-dead) {{\color{green!60!black}$\Xi_{2,2}$}};
        \draw[edge-k] (n1-3) -- (n2-k-from-dead);
        \node[node-k-dead] at (2, -2.4) {};
        \draw[edge-k-dashed] (n1-3) -- (2, -2.4);
        \node[node-k-dead] at (2, -1.6) {};
        \draw[edge-k-dashed] (n1-3) -- (2, -1.6);

        % ================= n=3 (核心耦合：左侧同步) =================
        % 蓝色最小值在 -3.0，绿色最小值在 -1.8
        % 定义相同的位移向量 Xi: {-1.2, -0.2, 0.5, 1.2}

        % --- 绿色最左侧的四分叉 (必须与上方蓝色四分叉完全平行) ---
        \node[node-k] (n3-k-1) at (3, -1.8) {}; % -0.6 + (-1.2)
        \node[node-k] (n3-k-2) at (3, -0.8) {}; % -0.6 + (-0.2)
        \node[node-k-dead] (n3-k-3) at (3, -0.1) {};
        \node[node-k-dead] (n3-k-4) at (3, 0.4) {};
        \draw[edge-k] (n2-k-min) -- (n3-k-1);
        \draw[edge-k] (n2-k-min) -- (n3-k-2);
        \draw[edge-k-dashed] (n2-k-min) -- (n3-k-3);
        \draw[edge-k-dashed] (n2-k-min) -- (n3-k-4);

        % --- 蓝色最左侧的四分叉 ---
        \node[node-N] (n3-N-1) at (3, -4.2) {}; % -3.0 + (-1.2)
        \node[node-N] (n3-N-2) at (3, -3.2) {}; % -3.0 + (-0.2)
        \node[node-N-dead] (n3-N-3) at (3, -2.5) {};
        \node[node-N-dead] (n3-N-4) at (3, -2.0) {};
        \draw[edge-N] (n2-N-min) -- (n3-N-1);
        \draw[edge-N] (n2-N-min) -- (n3-N-2);
        \draw[edge-N-dashed] (n2-N-min) -- (n3-N-3);
        \draw[edge-N-dashed] (n2-N-min) -- (n3-N-4);
        \node[below] at (n2-N-min) {{\color{blue!70!black}$\Xi_{1,2}$}};

        % 绿色其它支系
        \node[node-k] (n3-k-other) at (3, 1.3) {};
        \draw[edge-k] (n2-k-from-dead) -- (n3-k-other);
        \draw[edge-N-dashed] (2, -2.2) -- (3, -1.1);
        \node[node-N-dead] at (3, -1.1) {};
        \node[right] at (2, -2.2) {{\color{blue!70!black}$\Xi_{2,2}$}};

        % ================= 标注 =================
        % 最小值对比
        \draw[blue!70!black, <-, thick] (3.05, -4.2) -- (3.15, -4.2) node[right] {$\mathcal{M}^2_3$};
        \draw[green!60!black, <-, thick] (3.05, -1.8) -- (3.15, -1.8) node[right] {$\mathcal{M}^{\mathrm{k}}_3$};

        % % 耦合箭头：显示位移的一致性
        % \draw[<->, gray, thin, bend left=20] (2.1, -3.0) to node[midway, right, font=\tiny] {$\Xi_{1,n}$} (2.1, -1.8);
        % \draw[stealth-stealth, orange, dashed] (2.9, -4.2) -- (2.9, -3.0) node[midway, left, font=\tiny] {Parallel $\Xi$};

    \end{tikzpicture}
    \caption{\bilingual{Illustration of a \texorpdfstring{$2$}{2}-BRW (blue) and a killed BRW (green) coupling}{一个\texorpdfstring{$2$-BRW}{2-BRW} (蓝色) 与一个经灭杀的\texorpdfstring{BRW}{BRW} (绿色) 的耦合示例图}}
\end{figure}
